\documentclass[10pt,aspectratio=169]{beamer}
\usepackage[utf8]{inputenc}
\usepackage[T1]{fontenc}
\usepackage[ngerman]{babel}
\usepackage{lmodern}
\usepackage{booktabs,tabularx,array}
\usepackage{csquotes}
\usepackage{tikz}
\usetikzlibrary{positioning,calc}

\usetheme{default}
\usecolortheme{beaver}
\setbeamertemplate{navigation symbols}{}
\setbeamertemplate{footline}{\hfill\scriptsize\insertframenumber/3\hspace*{4mm}\vspace*{2mm}}
\setbeamerfont{frametitle}{size=\large,series=\bfseries}
\newcommand{\code}[1]{\texttt{#1}}
\newcolumntype{L}{>{\raggedright\arraybackslash}X}
\renewcommand{\arraystretch}{1.1}

\begin{document}

%---------------------------------------------------------------------
\begin{frame}{1 \,|\, Frankfurt im Graphen: 10 Fragen, 8 Tools}
\vspace{-2mm}
\small Gruppe 2 · RI XIII H.\,4 (Stadtarchiv Frankfurt) · \textbf{1.054 Regesten, 1440--1493} · gemeinsames Modell RI-alles
\vspace{1mm}

\scriptsize
\begin{tabularx}{\textwidth}{@{}rLp{5.3cm}p{3.1cm}@{}}
\toprule
\# & Frage der Lehrkraft & Modell (Knoten / Annotation / Rolle) & Tool\\\midrule
1 & Welche Urkunden stellte F.\ 1442 in Frankfurt aus? & \code{ri:placeOfIssue}, Datum & \code{search\_regesta}\\
2 & Von wo aus schrieb Friedrich an Frankfurt? & \code{ri:placeOfIssue} $\to$ Place & \code{itinerary}\\
3 & Welche Privilegien erhielt die Stadt? & \textbf{\code{ri:legalAct}} \code{privilege\_*} & \code{legal\_acts}\\
4 & Welche Steuern zahlte Frankfurt? & Summary + \textbf{\code{receipt}, \code{year\_ref}} & \code{legal\_acts}\\
5 & Was regelte der Landfrieden 1442? & Regesta n.\,41, Summary & \code{get\_regest}\\
6 & Welche Kanzleibeamten wirkten mit? & \code{nennung} + \textbf{\code{relator}/\code{chancery}} & \code{regesta\_by\_entity}\\
7 & Welche Befehle an den Rat -- gegen wen? & \textbf{\code{mandate}} + \textbf{Rolle \code{affected}} & \code{legal\_acts}\\
8 & Wie entwickelte sich der Kontakt? & \code{startDate}/\code{endDate} & \code{count\_regesta}\\
9 & Wer wurde zusammen mit Frankfurt angeschrieben? & \code{ri:appearsIn \{role:'empfaenger'\}} & \code{co\_recipients}\\
10 & Welche Personen begünstigte der König? & \textbf{Rolle \code{beneficiary}} & \code{regesta\_by\_entity}\\
\bottomrule
\end{tabularx}

\vspace{2mm}
\small
\begin{columns}[T]
\column{0.45\textwidth}
\textbf{5} Basis · \textbf{2} genauer mit Ergänzung · \textbf{3} nur mit Ergänzung (\textbf{fett})
\column{0.52\textwidth}
\textbf{Grundsatz:} Kein freies Cypher -- das LLM ruft \textbf{feste, getestete MCP-Tools} auf. Jede Antwort nennt Regestennummer und Link.
\end{columns}
\end{frame}

%---------------------------------------------------------------------
\begin{frame}{2 \,|\, Fallstricke und unsere Ergänzung}
\begin{columns}[T]
\column{0.47\textwidth}
\textbf{Fallstricke -- fest in den Tools}\\[1mm]
\footnotesize
\textbf{Chmel-Dubletten:} 100 Regesten verlinken Chmel (z.\,T.\ mit anderem Datum).
$\to$ Heft zählt, Chmel nur als Hinweis.\\[1.5mm]
\textbf{Datum als Spanne:} 483 von 1.054 unsicher; \enquote{[Vor 1441 Mai 26]} steht als exakter Tag.
$\to$ Überlappung \code{start <= Ende AND end >= Anfang}, Originaldatum + \enquote{erschlossen} ausgeben.\\[1.5mm]
\textbf{Fehlende Orte:} 442 (42\,\%) ohne Ort -- hängen am \textbf{Schein-Ort \enquote{-}}, nicht NULL.
$\to$ ausschließen, Lücke immer nennen, \emph{kein} Raten per LLM.\\[1.5mm]
\textbf{Außerdem:} \code{count(DISTINCT r)} (doppelte Annotationen), \enquote{Frankfurt} hat 4 Lesarten.

\column{0.5\textwidth}
\textbf{Ergänzung per LLM: Rollen + Rechtsakte}\\[1mm]
\footnotesize
n.\,32: \emph{Befehl an Schultheiß Hermann von Hoch-Weisel, Katharina von Ingelheim zu schützen}\\[1mm]
\begin{tabular}{@{}lll@{}}
\toprule
& Register & + LLM\\\midrule
Hoch-Weisel & \code{empfaenger} & \code{commissioned}\\
Katharina & \code{nennung} & \code{beneficiary}\\
Regest & -- & \code{ri:legalAct mandate}\\
\bottomrule
\end{tabular}\\[1.5mm]
Gleiches Muster wie RI-alles: \code{Annotation \{type, role|act,}\\\code{evidence, source, status\}}\\[1mm]
\textbf{Herkunft:} \code{register} · \code{llm} (Modell, Prompt, Belegstelle, \code{unreviewed}/\code{verified}) · \code{manual}\\[1mm]
\textbf{1.101 Rechtsakte, 1.456 Rollen}, 0 Fehler in der Autoprüfung; Goldstandard (50 Regesten) läuft.
\end{columns}
\end{frame}

%---------------------------------------------------------------------
\begin{frame}{3 \,|\, Architektur}
\begin{center}
\resizebox{!}{0.72\textheight}{% Architektur-Skizze (TikZ), wird von plan.tex und pitch.tex eingebunden
\begin{tikzpicture}[
  font=\sffamily\footnotesize,
  box/.style={draw, rounded corners=2pt, align=center, inner sep=4pt, minimum height=9mm},
  arr/.style={<->, >=stealth, thick},
  lab/.style={font=\sffamily\scriptsize, fill=white, inner sep=1pt},
  node distance=6mm
]
  % Frontend
  \node[box, fill=blue!10, text width=30mm] (ui) {\textbf{Browseroberfläche}\\(React)\\[2pt]
    \scriptsize Chat · Quellenbelege\\\scriptsize Karte · Zeitleiste\\\scriptsize Netzwerk · Steckbrief};

  % Backend
  \node[box, fill=orange!15, text width=40mm, right=17mm of ui] (bot) {\textbf{Chatbot}\\Python / FastAPI\\Agentic Harness\\[2pt]
    \scriptsize Systemprompt · Tool-Aufrufe prüfen\\\scriptsize Belege + Visualisierungsdaten};
  \node[box, fill=orange!30, text width=40mm, below=4mm of bot, minimum height=6mm] (mcpc) {\textbf{MCP-Client}};

  % LLM
  \node[box, fill=gray!20, text width=30mm, right=15mm of bot] (llm) {\textbf{LLM}\\(konfigurierbar)\\[2pt]
    \scriptsize THM-KI-Dienst oder lokal\\\scriptsize (llama.cpp / vLLM)};

  % MCP servers
  \node[box, fill=violet!12, text width=44mm, below=11mm of mcpc, xshift=-12mm] (rimcp) {\textbf{RI-MCP-Server} (unser)\\[2pt]
    \scriptsize\ttfamily search\_regesta · get\_regest\\\scriptsize\ttfamily find\_entity · itinerary\\
    \scriptsize\ttfamily regesta\_by\_entity · count\_regesta\\\scriptsize\ttfamily co\_recipients · legal\_acts\\[1pt]
    \scriptsize\sffamily feste, getestete Cypher-Abfragen};
  \node[box, fill=green!12, text width=28mm, right=6mm of rimcp] (wd) {\textbf{Wikidata-}\\\textbf{MCP-Server}\\[2pt]\scriptsize Steckbriefe\\\scriptsize (externe Quelle)};

  % DB
  \node[box, fill=violet!25, text width=60mm, below=9mm of rimcp] (db) {\textbf{Neo4j: RI-alles} (\texttt{riallesatag})\\
    \scriptsize + unsere Ergänzung: \texttt{ri:legalAct}, neue Rollen in \texttt{ri:appearsIn}};
  \node[box, dashed, text width=30mm, left=8mm of db] (jlu) {\scriptsize RI-alles auf jlu-buster\\\scriptsize\texttt{:9011} (schreibgeschützt)};
  \node[box, dashed, text width=34mm, right=8mm of db] (pipe) {\scriptsize LLM-Extraktion (offline)\\\scriptsize Rollen + Rechtsakte\\\scriptsize \texttt{source:'llm'}};

  \draw[arr] (ui) -- node[lab, above=1pt] {REST/JSON} (bot);
  \draw[arr] (bot) -- node[lab, above=1pt] {API} (llm);
  \draw[arr] (bot) -- (mcpc);
  \draw[arr] (mcpc.south) ++(-12mm,0) -- node[lab, right] {MCP} (rimcp.north);
  \draw[arr] (mcpc.south) ++(14mm,0) |- ++(0,-4mm) -| node[lab, pos=0.75, right] {MCP} (wd.north);
  \draw[->, >=stealth, thick] (rimcp) -- node[lab, right] {Bolt, nur lesend} (db);
  \draw[->, >=stealth, dashed] (jlu) -- node[lab, above] {Kopie} (db);
  \draw[->, >=stealth, dashed] (pipe) -- node[lab, above] {Import} (db);
\end{tikzpicture}
}
\end{center}
\vspace{-1mm}
\centering\large\textbf{Das LLM formuliert -- die Tools rechnen und belegen.}
\end{frame}

\end{document}
